\chapter{Instructions Communes}

\section{Structure du D\'ep\^ot}

Votre d\'ep\^ot Git doit respecter la convention de nommage suivante pour
tous les modules soumis. L'outil d'\'evaluation (\texttt{pendragon})
se base sur cette structure pour localiser et tester votre travail.

\begin{console}
<login>/
+-- <MODULE-ID>-<slug>/         # ex. M-BLS-001-business-logic/
|   +-- <turnindir>/            # ex. ex00-pricing-manipulation/
|   |   +-- <file1>             # ex. poc-pricing.py
|   |   +-- <file2>             # ex. report.md
|   |   +-- ...
|   +-- <turnindir>/            # ex. ex01-race-conditions/
|   |   +-- ...
|   +-- ...
+-- ...
\end{console}

\begin{itemize}
  \item \textbf{R\'epertoire du module} : le nom complet du r\'epertoire tel
        qu'affich\'e dans votre curriculum (ex. \texttt{M-BLS-001-business-logic}).
  \item \textbf{R\'epertoire de rendu} : la valeur de \texttt{\textbackslash turnindir}
        fournie dans chaque exercice (ex. \texttt{ex00-pricing-manipulation}).
  \item \textbf{Fichiers de l'exercice} : exactement les fichiers list\'es dans
        \texttt{\textbackslash exfiles}. Ni fichiers suppl\'ementaires, ni fichiers manquants.
\end{itemize}

\section{Qu'est-ce qu'une Preuve de Concept Valide}

Une PoC doit \^etre :

\begin{enumerate}
  \item \textbf{Reproductible}. N'importe qui avec le m\^eme environnement peut
        l'ex\'ecuter et observer la vuln\'erabilit\'e. Les chemins cod\'es en dur,
        les adresses IP absolues et les jetons expir\'es ne sont pas acceptables.
  \item \textbf{Document\'ee}. Chaque \'etape est expliqu\'ee. L'\'evaluteur ne
        devrait pas avoir \`a deviner ce que fait le script.
  \item \textbf{Cibl\'ee}. L'exercice sp\'ecifie une cible autoris\'ee
        (lab, maquette, binaire local ou plateforme d\'esign\'ee). Ne testez
        jamais les PoC sur des syst\`emes de production qui ne vous appartiennent pas.
  \item \textbf{Minimale}. La PoC d\'emontre \emph{uniquement} la vuln\'erabilit\'e
        en question. Les charges utiles non li\'ees et le bruit sont p\'enalis\'es.
\end{enumerate}

\begin{warning}
  Une PoC qui modifie silencieusement l'\'etat du syst\`eme, supprime des
  fichiers ou persiste des changements en dehors du p\'erim\`etre de l'exercice
  sera rejet\'ee. Si votre exploit n\'ecessite un nettoyage, documentez les
  \'etapes de nettoyage dans votre rapport.
\end{warning}

\section{Processus de Soumission}

\begin{enumerate}
  \item R\'ealisez les livrables de l'exercice sur votre machine locale.
  \item Committez et poussez vers votre d\'ep\^ot Git assign\'e.
  \item Le d\'ep\^ot est gel\'e apr\`es la date limite.
  \item Les pairs clonent votre d\'ep\^ot et \'evaluent en utilisant :
        \begin{itemize}
          \item \texttt{pendragon} pour l'ex\'ecution automatis\'ee des tests
                (l\`a o\`u le module fournit des configurations mouli).
          \item La grille manuelle pour la qualit\'e du rapport, le style de
                code et la m\'ethodologie.
        \end{itemize}
\end{enumerate}

\hint{Ex\'ecutez toujours \texttt{pendragon run <MODULE-ID> <ex\_id>
<path/to/student/repo>} avant votre commit final pour v\'erifier
que votre soumission passe les v\'erifications automatis\'ees.}

\subsection*{Test avec pendragon}

Si le module inclut une configuration mouli (fichiers JSON dans
\texttt{resources/mouli/}), vous pouvez tester votre soumission localement :

\begin{console}
pendragon run M-BLS-001 ex00 ~/student-repos/tedsig42
\end{console}

\texttt{pendragon} r\'esout le chemin de l'exercice automatiquement en
cherchant \texttt{<mod-fullname>/<ex-dirname>/} dans le r\'epertoire
\'etudiant donn\'e.

\subsection*{V\'erification de style avec caliban}

Pour les modules qui n\'ecessitent du code C, \texttt{caliban} v\'erifie
la conformit\'e du style :

\begin{console}
caliban check ~/student-repos/tedsig42/M-BLS-001-business-logic
\end{console}

\texttt{caliban} applique la Norme 42 (taille des fonctions, nommage
des variables, constructions interdites) adapt\'ee aux contextes de
recherche de vuln\'erabilit\'es. Ex\'ecutez-le avant la soumission ---
un \'echec aux v\'erifications de style peut invalider l'exercice.

\section{R\`egle des Fichiers Suppl\'ementaires Interdits}

\begin{itemize}
  \item Votre d\'ep\^ot doit contenir \textbf{uniquement} les fichiers
        list\'es dans \texttt{\textbackslash exfiles} pour chaque exercice.
  \item Fichiers interdits : \texttt{.DS\_Store}, \texttt{\_\_pycache\_\_},
        \texttt{.venv/}, \texttt{node\_modules/}, \texttt{.git/},
        binaires compil\'es (sauf explicitement list\'es dans \texttt{exfiles}),
        fichiers d'\'echange d'\'editeur, fichiers journaux et tout fichier
        non demand\'e.
  \item Utilisez \texttt{.gitignore} pour exclure automatiquement les
        fichiers ind\'esirables.
\end{itemize}

\section{Format du Rapport}

Lorsqu'un rapport (\texttt{report.md}) est requis, il doit inclure :

\begin{enumerate}
  \item \textbf{Classe de vuln\'erabilit\'e} (identifiant CWE si applicable).
  \item \textbf{Étapes de reproduction} --- num\'erot\'ees, pr\^etes
        \`a \^etre copi\'ees-coll\'ees.
  \item \textbf{Charges utiles} --- chaînes exactes, requ\^etes ou
        entr\'ees utilis\'ees.
  \item \textbf{Impact} --- ce qu'un attaquant peut accomplir.
  \item \textbf{Mitigation} --- comment le d\'eveloppeur devrait corriger.
  \item \textbf{Sources} --- toutes les ressources consult\'ees (voir
        Section~\ref{sec:cite-sources}).
\end{enumerate}

\subsection*{En-t\^ete de Rendu}

Chaque fichier source doit inclure l'en-t\^ete de rendu d\'efini dans
\path{../../../resources/latex/deliverable-header.tex}, adapt\'e
\`a la syntaxe de commentaire du langage utilis\'e.

\section{Cibles Autoris\'ees Uniquement}

\begin{warning}
  Vous pouvez \textbf{uniquement} tester les cibles explicitement list\'ees
  dans l'exercice (labs, serveurs locaux, maquettes, binaires d\'esign\'es).
  Tester des syst\`emes de production r\'eels, des applications tierces
  ou toute cible qui ne vous appartient pas ou pour laquelle vous n'avez
  pas d'autorisation \'ecrite est interdit et peut entraîner une
  disqualification imm\'ediate.
\end{warning}

L'\'ethique du bug bounty s'applique tout au long de ce curriculum :
\begin{itemize}
  \item Aucun test sur des programmes en production sans un accord de
        p\'erim\`etre valide.
  \item Aucun d\'eni de service sauf si l'exercice l'exige explicitement.
  \item Aucune exfiltration de donn\'ees au-del\`a de ce qui est
        n\'ecessaire pour d\'emontrer la vuln\'erabilit\'e.
  \item Si vous acc\'edez accidentellement \`a des donn\'ees r\'eelles
        d'utilisateurs, arr\^etez imm\'ediatement et documentez ce qui
        s'est pass\'e.
\end{itemize}

\section{Citez Toutes Vos Sources}
\label{sec:cite-sources}

Toute ressource consult\'ee pendant votre travail doit \^etre cit\'ee.
Cela inclut :

\begin{itemize}
  \item Articles acad\'emiques et billets de blog.
  \item Rapports publics de bug bounty (HackerOne, Bugcrowd, etc.).
  \item Documentation officielle (RFC, pages de manuel, documentation
        fournisseur).
  \item \textbf{Outils d'IA} : nom du mod\`ele, r\'esum\'e de la
        question, date et m\'ethode de v\'erification (voir le chapitre
        Instructions IA pour le format de citation exact).
  \item Extraits de code adapt\'es de projets open-source (avec
        attribution de licence si requis).
\end{itemize}

\subsection*{Format de citation}

\begin{minted}{markdown}
# R\'ef\'erences

- [1] PortSwigger, "Business Logic Flaws",
  https://portswigger.net/web-security/logic-flaws
  (consult\'e le 2025-03-15).
- [2] Rapport HackerOne #123456, "IDOR on /api/users",
  anonymis\'e, consult\'e via divulgation publique.
- [3] G\'en\'er\'e avec l'assistance de Claude 3.5 Sonnet
  le 2025-03-14. Question : "Write a Python script that...".
  V\'erifi\'e par : revue manuelle de chaque requ\^ete HTTP.
\end{minted}

\subsection*{Attribution en ligne}

Si un exercice vous demande explicitement d'utiliser une technique ou
un outil sp\'ecifique, vous n'avez pas besoin de citer l'exercice
lui-m\^eme comme source. Citez tout ce qui se trouve en dehors de
l'\'enonc\'e de l'exercice.
